\documentclass[11pt,a4paper]{article}
\usepackage[utf8]{inputenc}
\usepackage[T1]{fontenc}
\usepackage[brazilian]{babel}
\usepackage{geometry}
\geometry{top=1.8cm, bottom=1.8cm, left=1.8cm, right=1.8cm}
\usepackage{xcolor}
\usepackage{titlesec}
\usepackage{enumitem}
\usepackage{booktabs}
\usepackage{tabularx}
\usepackage{tcolorbox}
\usepackage{hyperref}
\usepackage{microtype}
\usepackage{newtxtext,newtxmath}

% Definição de Cores Institucionais IFES / Profissional
\definecolor{ifesgreen}{RGB}{24, 138, 63}
\definecolor{ifesred}{RGB}{200, 30, 30}
\definecolor{darkgray}{RGB}{50, 50, 50}
\definecolor{lightgray}{RGB}{245, 247, 248}
\definecolor{bordergray}{RGB}{210, 215, 220}
\definecolor{boxheader}{RGB}{235, 243, 238}

\hypersetup{
    colorlinks=true,
    linkcolor=ifesgreen,
    urlcolor=ifesgreen,
    citecolor=ifesgreen
}

% Formatação de Títulos
\titleformat{\section}{\color{ifesgreen}\large\bfseries}{\thesection}{1em}{}[\color{ifesgreen}\titlerule]
\titleformat{\subsection}{\color{darkgray}\normalsize\bfseries}{\thesubsection}{1em}{}

\tcbset{
    colback=lightgray,
    colframe=ifesgreen,
    coltitle=white,
    fonttitle=\bfseries\small,
    boxrule=0.8pt,
    arc=3pt,
    left=6pt, right=6pt, top=6pt, bottom=6pt
}

\begin{document}

% Cabeçalho
\begin{center}
    {\small\textbf{INSTITUTO FEDERAL DO ESPÍRITO SANTO -- IFES | CAMPUS SERRA}}\\[2pt]
    {\small\textbf{PROGRAMA DE PÓS-GRADUAÇÃO EM ENGENHARIA DE CONTROLE E AUTOMAÇÃO (PROPECAUT)}}\\[6pt]
    \rule{\textwidth}{1.2pt}\\[6pt]
    {\large\textbf{\color{ifesgreen}PROPOSTA DE REESTRUTURAÇÃO TEMÁTICA DA DISSERTAÇÃO}}\\[3pt]
    {\small\textbf{Alinhamento de Escopo, Foco em Controle e Definição de Fronteiras de Modelagem}}\\[4pt]
    \rule{\textwidth}{0.6pt}
\end{center}

\vspace{-6pt}
\begin{center}
    \begin{tabularx}{\textwidth}{lX}
        \textbf{Título Oficial:} & \textbf{Controle e Gerenciamento Térmico-Energético Computacional de um Modelo de Data Center Orbital} \\
        \textbf{Mestrando:} & Dafnhê Thor Castro de Martins Barros \\
        \textbf{Orientador:} & Prof. Dr. Flávio Barcelos Braz da Silva \\
        \textbf{Objetivo do Memo:} & Apresentar a estrutura comparativa (Como está $\times$ Como pretendemos que fique) para validação da orientação antes da reescrita dos Capítulos 3 e 4.
    \end{tabularx}
\end{center}

\vspace{4pt}

% Box com os 5 Pilares Centrais
\begin{tcolorbox}[colback=boxheader, colframe=ifesgreen, title=\color{ifesgreen}\textbf{DIRETRIZ CENTRAL: OS 5 PILARES DA DISSERTAÇÃO ENXUTA}]
\small
\begin{enumerate}[leftmargin=1.5em, itemsep=1pt, parsep=0pt]
    \item \textbf{Entrada de Energia:} Balanço macro de geração fotovoltaica AM0 e suporte do banco de baterias (BESS).
    \item \textbf{Geração de Calor:} Dinâmica elétrica e térmica de cargas de IA nos semicondutores ($P_{chip} = \alpha C V_{dd}^2 f_{clk} + P_{leak}$).
    \item \textbf{Física da Dissipação no Vácuo:} Condução estrutural interna e rejeição radiativa via Stefan-Boltzmann ($q \propto T_{rad}^4 - T_{space}^4$).
    \item \textbf{Modelo Preditivo (PINN):} Rede neural informada por física resolvendo o campo térmico em $<5\text{ ms}$ para o computador de bordo.
    \item \textbf{Controle Ativo (TEMS / DVFS):} Algoritmo em malha fechada modulando clock, tensão e tarefas para garantir $T_j \le 80^\circ\text{C}$.
\end{enumerate}
\end{tcolorbox}

\vspace{4pt}
\section{Esqueleto Comparativo Capítulo a Capítulo}

\subsection*{Capítulo 1: Introdução}
\begin{itemize}[leftmargin=1.2em, itemsep=2pt, parsep=0pt]
    \item \textbf{Como está hoje:} Contextualiza a crise hídrica/energética terrestre de IA e a solução orbital; formula o paradoxo térmico; lista 8 objetivos específicos (detalhando perdas magnéticas de ferrita e fadiga de soldas BGA).
    \item \textbf{Como pretendemos que fique:} \textbf{Mantém a motivação e enxuga os objetivos.} Foca os objetivos específicos na integração dos 5 pilares, eliminando temas periféricos de microeletrônica/mecânica.
    \item \textbf{Ajuste terminológico crucial:} Substituir a menção a ``Gêmeo Digital embarcado'' por ``Modelo Substituto Preditivo baseado em Física (\textit{Physics-Informed Surrogate Model})''.
\end{itemize}

\subsection*{Capítulo 2: Fundamentação Teórica e Estado da Arte}
\begin{itemize}[leftmargin=1.2em, itemsep=2pt, parsep=0pt]
    \item \textbf{Como está hoje:} Fundamenta o ambiente espacial, órbitas (SSO, LEO, GEO), lei de Stefan-Boltzmann, subsistema elétrico, estimador de bateria AUKF e fundamentos de PINN. (A subseção de custos do foguete Falcon 9 já foi removida e substituída pela justificativa do nó de 50~kW).
    \item \textbf{Como pretendemos que fique:} \textbf{Reequilibrar os pesos temáticos.} Reduzir o excesso de detalhes eletroquímicos da bateria e expandir a Seção do TEMS (que hoje tem apenas 2 parágrafos), evidenciando a formulação de controle preditivo e modulação DVFS.
\end{itemize}

\subsection*{Capítulo 3: Modelagem do Subsistema Elétrico e de Cargas}
\begin{itemize}[leftmargin=1.2em, itemsep=2pt, parsep=0pt]
    \item \textbf{Como está hoje:} Extensamente fragmentado com tópicos de eletrônica analógica pesada: perdas magnéticas em ferrita MnZn (Steinmetz), estabilidade de impedância de Middlebrook, corte de arco em chaves SSPC em microssegundos e filtro AUKF matricial de 2ª ordem.
    \item \textbf{Como pretendemos que fique (Enxuto - Entrada de Energia e Consumo de IA):}
    \begin{enumerate}[label=\alph*., itemsep=1pt, parsep=0pt]
        \item \textbf{Balanço de Potência do Nó (50 kW):} Equacionamento macro da distribuição HVDC ($P_{gerada} = P_{TI} + P_{bat} + P_{perdas}$) com rendimento de conversores GaN ($\eta \approx 98\%$).
        \item \textbf{Modelo da Fonte Primária (Arranjo AM0):} Equacionamento consolidado das células multijunção sob irradiância orbital e temperatura.
        \item \textbf{Balanço Eletroquímico de Suporte (BESS):} Dinâmica de Estado de Carga ($SoC$), rendimento coulômbico e restrições de ciclagem ($DoD \le 20\text{--}30\%$).
        \item \textbf{Modelagem das Cargas de IA (Injeção de Calor):} Relação analítica do consumo dos processadores neurais ($P_{chip} = \alpha C V_{dd}^2 f_{clk} + V_{dd} I_{leak}$), conectando o processamento computacional à fonte de aquecimento.
    \end{enumerate}
    \item \textbf{O que sai:} Perdas de Steinmetz em núcleos de indutores, arcos de SSPC e estabilidade de Middlebrook (pertencentes a projetos de fontes, não ao controle térmico-energético do satélite).
\end{itemize}

\vspace{4pt}

\subsection*{Capítulo 4: Modelagem Térmica, Preditor PINN e Controle Ativo TEMS}
\begin{itemize}[leftmargin=1.2em, itemsep=2pt, parsep=0pt]
    \item \textbf{Como está hoje:} Apresenta fluxos orbitais e rede térmica nodal; porém, contém desvios disciplinares severos (fadiga termomecânica de esferas de solda BGA por Coffin-Manson e matriz FMECA industrial de 17 componentes), e a PINN ainda não está acoplada na malha fechada.
    \item \textbf{Como pretendemos que fique (O Coração da Dissertação de Controle):}
    \begin{enumerate}[label=\alph*., itemsep=1pt, parsep=0pt]
        \item \textbf{Física da Dissipação no Vácuo:} Condução transiente no chassi/dissipador acoplada à radiação de Stefan-Boltzmann para o espaço profundo ($3\text{ K}$).
        \item \textbf{Modelo Preditivo PINN (DeepXDE/PyTorch):} Treinamento da rede neural informando a equação do calor e a condição de contorno radiativa na função de perda; validação da inferência térmica contínua em $< 5\text{ ms}$ (viável para o computador de bordo OBC).
        \item \textbf{Sistema de Gerenciamento Térmico-Energético Ativo (TEMS):} Algoritmo de controle preditivo em malha fechada utilizando as estimativas da PINN para modular DVFS ($f_{clk}, V_{core}$) e escalonar tarefas de IA, impondo a restrição rígida $T_j \le 80^\circ\text{C}$.
        \item \textbf{Resultados e Validação Comparativa:} Simulação de 3 cenários: (1) Sem controle (superaquecimento $> 100^\circ\text{C}$), (2) Controle reativo convencional por corte abrupto (\textit{thermal throttling}), (3) Controle preditivo ativo TEMS+PINN (operação contínua e suave a $\le 80^\circ\text{C}$).
    \end{enumerate}
    \item \textbf{O que sai:} Fadiga mecânica de solda BGA e análise FMECA (desvios mecânicos/industriais que não agregam ao modelo de controle).
\end{itemize}

\subsection*{Capítulo 5: Conclusões e Trabalhos Futuros}
\begin{itemize}[leftmargin=1.2em, itemsep=2pt, parsep=0pt]
    \item \textbf{Como está hoje:} Síntese genérica dos subsistemas e menção a testes laboratoriais futuros.
    \item \textbf{Como pretendemos que fique:}
    \begin{enumerate}[label=\alph*., itemsep=1pt, parsep=0pt]
        \item \textbf{Conclusões:} Quantificação da eficácia do TEMS+PINN na contenção térmica com máxima vazão de IA.
        \item \textbf{Limitações Epistêmicas:} Fronteiras da modelagem estritamente computacional (\textit{in silico}).
        \item \textbf{Trabalhos Futuros (Posicionamento do Gêmeo Digital):} Desenvolvimento de \textbf{Gêmeo Digital Operacional} com calibração contínua por telemetria orbital real em tempo real; ensaios em câmara de termovácuo (TVAC); e expansão para redes com múltiplos satélites em constelação.
    \end{enumerate}
\end{itemize}

\vspace{4pt}
\section{Tabela Síntese das Modificações Propostas}

\begin{table}[h!]
\centering
\small
\begin{tabularx}{\textwidth}{p{3.8cm}Xp{4.2cm}}
    \toprule
    \textbf{Capítulo / Seção} & \textbf{Conteúdo Proposto (Foco em Controle e Automação)} & \textbf{O que Foi Eliminado (Desvio Temático)} \\
    \midrule
    \textbf{Cap. 1: Introdução} & Foco na inviabilidade de FEM/CFD em tempo real no OBC e necessidade de preditor rápido. & Menção a Gêmeo Digital embarcado (ajustado para Surrogate Model). \\
    \addlinespace
    \textbf{Cap. 2: Fundamentação} & Física de radiação, ambiente AM0, nó modular de 50 kW, arquitetura TEMS e bases da PINN. & Análise econômica e custos de voo do Falcon 9 / coifa 5,2 m (removidos). \\
    \addlinespace
    \textbf{Cap. 3: Elétrico e Cargas} & Balanço macro de potência (50 kW), arranjo solar, dinâmica de SoC do BESS e injeção térmica do chip de IA. & Perdas em ferrita MnZn (Steinmetz), arcos de chaves SSPC e estabilidade Middlebrook. \\
    \addlinespace
    \textbf{Cap. 4: Térmico e Controle} & Condução e radiação no vácuo, preditor PINN em tempo real ($<5\text{ ms}$) e controle TEMS/DVFS ($T_j \le 80^\circ\text{C}$). & Fadiga termomecânica de solda BGA (Coffin-Manson) e matriz FMECA 17 modos. \\
    \addlinespace
    \textbf{Cap. 5: Conclusões} & Fechamento quantitativo dos 5 pilares e posicionamento formal de Gêmeo Digital como trabalho futuro. & Promessas de validações experimentais fora do escopo da dissertação. \\
    \bottomrule
\end{tabularx}
\end{table}

\vspace{6pt}
\begin{tcolorbox}[colback=white, colframe=darkgray, title=\color{darkgray}\textbf{RESUMO DA ESTRATÉGIA PARA APRESENTAÇÃO AO ORIENTADOR}]
\small
\textbf{Mensagem-chave:} \textit{``Professor, identificamos que a dissertação continha trechos que fugiam da área de Controle e Automação (análise de custos de foguete, fadiga mecânica de soldas e física profunda de baterias). Estamos blindando a dissertação em torno de 5 pilares focados: Geração Solar $\to$ Consumo/Calor de IA $\to$ Dissipação no Vácuo $\to$ Preditor PINN em milissegundos $\to$ Controle Ativo TEMS/DVFS em malha fechada ($T_j \le 80^\circ\text{C}$). O conceito de Gêmeo Digital completo com telemetria fica posicionado como evolução futura no Capítulo 5.''}
\end{tcolorbox}

\end{document}
